%-----------------------------------------------------------------------
% Appendix A: User Manual
%-----------------------------------------------------------------------

\section{Introduction}
\label{app:a:introduction}

This user manual provides instructions for accessing, configuring, and playing the HPC Sorting Serious Game. It is intended for students, educators, and anyone interested in learning about parallel computing through interactive gameplay.

The game is deployed as a \textbf{web application} (Godot Engine WebAssembly export) and runs directly in a modern browser. No installation is required for the primary deployment. Desktop (Windows) builds are also available for offline use or development.

%-----------------------------------------------------------------------
\section{System Requirements}
\label{app:a:system-requirements}

\subsection{Web Version (Primary)}

\begin{itemize}
    \item \textbf{Browser}: Chromium-based browser (Chrome, Edge, Brave) or Firefox, version 145 or later
    \item \textbf{WebGL}: WebGL 2.0 support required
    \item \textbf{JavaScript}: Must be enabled
    \item \textbf{RAM}: 2~GB minimum available to browser
    \item \textbf{Network}: Internet connection required for multiplayer mode; single-player works offline after initial load
    \item \textbf{Input}: Mouse or touchscreen
\end{itemize}

%-----------------------------------------------------------------------
\section{Accessing the Game}
\label{app:a:accessing}

\subsection{Web Version}

\begin{enumerate}
    \item Navigate to the hosted game URL in a supported browser
    \item Wait for the WebAssembly module to load (progress bar shown)
    \item The main menu appears once loading is complete
\end{enumerate}

\noindent\textbf{Note:} The web export is built with thread support disabled (\texttt{variant/thread\_support=false} in the Godot export preset), so COOP/COEP headers and \texttt{SharedArrayBuffer} are \emph{not} required. The game runs as a single-threaded WebAssembly application. Multiplayer communication uses HTTP/SSE, which does not depend on shared memory or cross-origin isolation.

%-----------------------------------------------------------------------
\section{Getting Started}
\label{app:a:getting-started}

\subsection{Main Menu}

On launch, the main menu presents the following options:
\begin{itemize}
    \item \textbf{Single Player}: Practice sorting alone (sequential execution baseline)
    \item \textbf{Multiplayer}: Join or host a cooperative multiplayer game (OpenMP shared-memory simulation)
    \item \textbf{Options}: Configure theme and display options
\end{itemize}

%-----------------------------------------------------------------------
\section{Single-Player Mode}
\label{app:a:singleplayer-mode}

\subsection{Starting a Single-Player Game}

\begin{enumerate}
    \item Select ``Single Player'' from the main menu
    \item Configure game parameters in the options dialog:
          \begin{itemize}
              \item \textbf{Number of cards}: configurable via spin box
              \item \textbf{Buffer slot count}: number of local buffer positions
              \item \textbf{Card value range}: upper bound for card values
          \end{itemize}
    \item Click ``Start'' to begin
\end{enumerate}

\subsection{Gameplay}

\paragraph*{Objective:}
Sort all cards in ascending order in the main container.

\paragraph*{Controls:}
\begin{itemize}
    \item \textbf{Drag and Drop}: Click and drag a card to reposition it within the container or move it to/from a buffer slot
    \item \textbf{Scroll}: Use the horizontal scroll bar or mouse wheel to navigate through cards when they exceed viewport width
\end{itemize}

\paragraph*{Using Buffer Slots:}
\begin{itemize}
    \item Buffer slots appear below the main card container
    \item Drag cards into buffer slots for temporary storage during sorting
    \item Buffer slots simulate thread-local storage in parallel computing
    \item Cards in buffers must form a contiguous subsequence of the sorted target before the game can complete
\end{itemize}

\paragraph*{Winning:}
\begin{itemize}
    \item The game completes when all cards are in correct ascending order in the main container
    \item Completion time is displayed
\end{itemize}

%-----------------------------------------------------------------------
\section{Multiplayer Mode}
\label{app:a:multiplayer-mode}

\subsection{Creating a Game}

\begin{enumerate}
    \item Select ``Multiplayer'' from the main menu
    \item Select ``Create Game''
    \item Configure game settings:
          \begin{itemize}
              \item Number of cards
              \item Buffer slot count per player
              \item Barrier mode (first-to-reach or round-robin)
          \end{itemize}
    \item A lobby code is generated automatically
    \item Share this code with other players
    \item Wait for players to join the lobby
    \item Once all players are present, click ``Start Game''
\end{enumerate}

\subsection{Joining a Game}

\begin{enumerate}
    \item Select ``Multiplayer'' from the main menu
    \item Select ``Join Game''
    \item Enter the lobby code provided by the host
    \item Click ``Join''
    \item Wait in the lobby for the host to start the game
\end{enumerate}

\subsection{Multiplayer Gameplay}

\paragraph*{Shared Container (Shared Memory):}
\begin{itemize}
    \item All cards are visible in the shared container, simulating shared memory accessible by all threads
    \item Any player can drag cards within the shared container to reorder them
    \item Changes propagate to all connected players in real time
\end{itemize}

\paragraph*{Local Buffers (Partitioned Shared Memory):}
\begin{itemize}
    \item Each player has local buffer slots visible only to themselves
    \item Cards moved to a player's buffer disappear from other players' views
    \item When a card is returned to the shared container, it becomes visible to all again
    \item This simulates the distinction between shared and private memory in OpenMP
\end{itemize}

\paragraph*{Barrier Synchronization:}
\begin{itemize}
    \item Any player can initiate a barrier (analogous to \texttt{\#pragma omp barrier})
    \item All players must reach the barrier before proceeding
    \item Once all players are at the barrier, a designated ``main thread'' (one player) can perform coordinated operations while others are blocked
    \item Main thread selection is either first-to-reach or round-robin, depending on game settings
    \item This teaches students about synchronization primitives in parallel computing
\end{itemize}

\paragraph*{Winning Together:}
\begin{itemize}
    \item The team wins when all cards are sorted correctly in the shared container
    \item All players must have returned their buffer cards to the container
\end{itemize}

%-----------------------------------------------------------------------
\section{Settings}
\label{app:a:settings}

\begin{itemize}
    \item \textbf{Theme}: Toggle between light and dark mode
    \item \textbf{Signaling Server URL}: Configure a custom signaling server address for multiplayer (advanced; defaults to the hosted server)
\end{itemize}

%-----------------------------------------------------------------------
\section{Tips and Strategies}
\label{app:a:tips-strategies}

\subsection{Sorting Strategies}

\begin{enumerate}
    \item \textbf{Divide and Conquer}: Use buffer slots to sort smaller subsets of cards, then merge them back
    \item \textbf{Find Sorted Runs}: Identify cards that are already in order and build around them
    \item \textbf{Minimize Moves}: Think before moving---unnecessary rearrangements waste time
\end{enumerate}

\subsection{Multiplayer Coordination}

\begin{enumerate}
    \item \textbf{Divide Responsibility}: Each player takes a value range (e.g., player~1 handles cards 1--25, player~2 handles 26--50)
    \item \textbf{Use Barriers Strategically}: Initiate a barrier when all players have sorted their portions and need to merge
    \item \textbf{Communicate Externally}: Use voice chat or messaging to coordinate since the game does not include built-in chat
    \item \textbf{Main Thread Awareness}: When you are the main thread during a barrier, merge sorted chunks efficiently; other players must wait
\end{enumerate}

%-----------------------------------------------------------------------
\section{Troubleshooting}
\label{app:a:troubleshooting}

\subsection{Common Issues and Solutions}

\paragraph*{Game does not load in browser:}
\begin{itemize}
    \item Verify that your browser supports WebGL 2.0 (check \texttt{chrome://gpu} in Chrome)
    \item Verify that JavaScript is enabled and WebGL 2.0 is supported (check \texttt{chrome://gpu} in Chrome)
    \item Try clearing browser cache and reloading
    \item Try a different Chromium-based browser
\end{itemize}

\paragraph*{Performance issues:}
\begin{itemize}
    \item Reduce the number of cards
    \item Close other browser tabs to free memory
    \item Ensure hardware acceleration is enabled in browser settings
\end{itemize}

\paragraph*{Cannot connect to multiplayer:}
\begin{itemize}
    \item Verify internet connection is active
    \item Check that both players are using the same game version
    \item Ensure the signaling server is reachable (default server requires internet access)
    \item Re-enter the lobby code carefully
    \item If using a custom signaling server, verify it is running and accessible
\end{itemize}

\paragraph*{Multiplayer desynchronization:}
\begin{itemize}
    \item This can occur if network packets are lost or arrive out of order
    \item The host can use the ``sync state'' mechanism to push authoritative state to all clients
    \item As a last resort, restart the game session
\end{itemize}

%-----------------------------------------------------------------------
\section{Educational Use}
\label{app:a:educational-use}

\subsection{For Students}

\begin{itemize}
    \item Play single-player mode first to understand the card sorting mechanics
    \item Observe how buffer slots help organize your sorting work (thread-local storage)
    \item In multiplayer, notice how coordination becomes harder with more players (communication overhead)
    \item Pay attention to the barrier mechanic---it directly mirrors \texttt{\#pragma omp barrier}
    \item Reflect on how game concepts relate to OpenMP parallel programming
\end{itemize}

\subsection{For Instructors}

\begin{itemize}
    \item Use as a pre-lecture activity to introduce HPC concepts before formal instruction
    \item The web deployment requires no installation---students can play immediately via a shared URL
    \item Facilitate post-gameplay discussion connecting game mechanics to parallel programming:
          \begin{itemize}
              \item Shared container $\leftrightarrow$ shared memory
              \item local buffers $\leftrightarrow$ thread-local / private variables
              \item Barrier mechanic $\leftrightarrow$ \texttt{\#pragma omp barrier}
              \item Multiple players $\leftrightarrow$ multiple threads
              \item Single-player $\leftrightarrow$ sequential execution baseline
          \end{itemize}
    \item Use completion metrics to discuss relative effort of parallel coordination vs.\ sequential sorting
\end{itemize}

%-----------------------------------------------------------------------
\section{Support and Resources}
\label{app:a:support}

\begin{itemize}
    \item \textbf{Game Repository}: \url{https://github.com/Siponek/hpc-sorting-serious-game}
    \item \textbf{Thesis Repository}: \url{https://github.com/Siponek/HPC-thesis}
    \item \textbf{Issue Tracker}: Report bugs and request features via GitHub Issues
    \item \textbf{License}: MIT license---contributions are welcome
\end{itemize}

%-----------------------------------------------------------------------
